% ═══════════════════════════════════════════════════════════════
% ÜBUNG: Exotische Oxide – Formeln und Reaktionsgleichungen
% Chemie Klasse 8 | Wiederholung vor Reduktion
% ═══════════════════════════════════════════════════════════════

\documentclass[11pt, a4paper]{article}

\usepackage[utf8]{inputenc}
\usepackage[T1]{fontenc}
\usepackage[ngerman]{babel}

\usepackage{helvet}
\renewcommand{\familydefault}{\sfdefault}

\usepackage{microtype}
\usepackage[margin=2cm, top=2cm, bottom=1.5cm]{geometry}
\usepackage{enumitem}
\usepackage{tabularx}
\usepackage{booktabs}
\usepackage{array}

\usepackage{fancyhdr}
\usepackage{lastpage}

\usepackage{xcolor}
\usepackage{tcolorbox}
\tcbuselibrary{skins}

% ═══════════════════════════════════════════════════════════════
% FARBEN
% ═══════════════════════════════════════════════════════════════

\definecolor{chemie}{HTML}{2a7a4b}
\definecolor{hellgrau}{HTML}{F5F5F5}
\definecolor{mittelgrau}{HTML}{888888}
\definecolor{dunkelgrau}{HTML}{444444}

\colorlet{fachfarbe}{chemie}

% ═══════════════════════════════════════════════════════════════
% KOPF- UND FUSSZEILE
% ═══════════════════════════════════════════════════════════════

\pagestyle{fancy}
\fancyhf{}
\renewcommand{\headrulewidth}{0.4pt}
\renewcommand{\footrulewidth}{0pt}
\lhead{\small Chemie Klasse 8 | Reaktionsgleichungen}
\rhead{\small Name: \underline{\hspace{3cm}}}
\cfoot{\small\textcolor{mittelgrau}{Seite \thepage{} von \pageref{LastPage}}}

% ═══════════════════════════════════════════════════════════════
% BOXEN
% ═══════════════════════════════════════════════════════════════

\newtcolorbox{infobox}[1][]{
    colback=hellgrau,
    colframe=hellgrau,
    boxrule=0pt,
    arc=0pt,
    left=8pt, right=8pt, top=6pt, bottom=6pt,
    borderline west={3pt}{0pt}{fachfarbe},
    #1
}

\newtcolorbox{merkbox}{
    colback=hellgrau,
    colframe=hellgrau,
    boxrule=0pt,
    arc=0pt,
    left=8pt, right=8pt, top=6pt, bottom=6pt,
    borderline north={2pt}{0pt}{fachfarbe}
}

% ═══════════════════════════════════════════════════════════════
% BEFEHLE
% ═══════════════════════════════════════════════════════════════

\newcommand{\luecke}[1]{\underline{\hspace{#1}}}

\newcommand{\aufgabe}[2]{%
    \vspace{8pt}%
    \noindent{\large\bfseries\textcolor{fachfarbe}{Aufgabe #1:} #2}%
    \vspace{6pt}%
    \nopagebreak[4]%
}

\setlength{\parindent}{0pt}
\setlength{\parskip}{4pt}

\newcolumntype{L}[1]{>{\raggedright\arraybackslash}p{#1}}
\newcolumntype{C}[1]{>{\centering\arraybackslash}p{#1}}

% ═══════════════════════════════════════════════════════════════
% DOKUMENT
% ═══════════════════════════════════════════════════════════════

\begin{document}

% ═══════════════════════════════════════════════════════════════
% HEADER
% ═══════════════════════════════════════════════════════════════

\begin{center}
    {\Large\bfseries\textcolor{fachfarbe}{Übung: Exotische Oxide}}\\[0.3em]
    {\normalsize Oxidformeln und Reaktionsgleichungen aufstellen}
\end{center}
\vspace{-0.3em}
\hrule
\vspace{0.8em}

% ═══════════════════════════════════════════════════════════════
% ERINNERUNG
% ═══════════════════════════════════════════════════════════════

\begin{merkbox}
\textbf{Erinnerung – Kreuzregel:}\\[0.3em]
Die Wertigkeit des einen Elements wird zum Index des anderen.\\
Beispiel: Aluminium (III) + Sauerstoff (II) $\rightarrow$ Al$_2$O$_3$
\hfill
\begin{tabular}{|c|c|}
\hline
\textbf{O} & \textbf{II} \\
\hline
\end{tabular}
\end{merkbox}

\vspace{0.5em}

% ═══════════════════════════════════════════════════════════════
% AUFGABE A
% ═══════════════════════════════════════════════════════════════

\aufgabe{A}{Oxidformeln aufstellen}

Stelle mithilfe der Kreuzregel die Formel des jeweiligen Oxids auf.

\vspace{0.5em}

\begin{center}
\renewcommand{\arraystretch}{1.8}
\begin{tabular}{|C{0.5cm}|L{4cm}|C{2.5cm}|C{3cm}|L{4cm}|}
\hline
\textbf{\#} & \textbf{Name} & \textbf{Wertigkeit} & \textbf{Formel} & \textbf{Verwendung} \\
\hline
1 & Bariumoxid & Ba (II) & & Feuerwerk \\
\hline
2 & Natriumoxid & Na (I) & & Glasherstellung \\
\hline
3 & Chromoxid & Cr (III) & & Grünes Pigment \\
\hline
4 & Titandioxid & Ti (IV) & & Sonnencreme, Farbe \\
\hline
5 & Manganoxid & Mn (IV) & & Batterien \\
\hline
6 & Zinnoxid & Sn (IV) & & Keramik, Glasuren \\
\hline
7 & Silberoxid & Ag (I) & & Angelaufenes Silber \\
\hline
\end{tabular}
\end{center}

\vspace{1em}

% ═══════════════════════════════════════════════════════════════
% AUFGABE B
% ═══════════════════════════════════════════════════════════════

\aufgabe{B}{Reaktionsgleichungen aufstellen und ausgleichen}

Formuliere die vollständige Reaktionsgleichung für die Oxidation. Gleiche aus!

\vspace{0.5em}

\begin{center}
\renewcommand{\arraystretch}{2.2}
\begin{tabular}{|C{0.5cm}|L{3.5cm}|L{9.5cm}|}
\hline
\textbf{\#} & \textbf{Wortgleichung} & \textbf{Formelgleichung (ausgeglichen)} \\
\hline
1 & Barium + Sauerstoff & \luecke{3cm} + \luecke{1.5cm} $\rightarrow$ \luecke{3cm} \\
\hline
2 & Natrium + Sauerstoff & \luecke{3cm} + \luecke{1.5cm} $\rightarrow$ \luecke{3cm} \\
\hline
3 & Chrom + Sauerstoff & \luecke{3cm} + \luecke{1.5cm} $\rightarrow$ \luecke{3cm} \\
\hline
4 & Titan + Sauerstoff & \luecke{3cm} + \luecke{1.5cm} $\rightarrow$ \luecke{3cm} \\
\hline
5 & Mangan + Sauerstoff & \luecke{3cm} + \luecke{1.5cm} $\rightarrow$ \luecke{3cm} \\
\hline
6 & Zinn + Sauerstoff & \luecke{3cm} + \luecke{1.5cm} $\rightarrow$ \luecke{3cm} \\
\hline
7 & Silber + Sauerstoff & \luecke{3cm} + \luecke{1.5cm} $\rightarrow$ \luecke{3cm} \\
\hline
\end{tabular}
\end{center}

\vspace{1em}

% ═══════════════════════════════════════════════════════════════
% ÜBERLEITUNG
% ═══════════════════════════════════════════════════════════════

\begin{infobox}
\textbf{Ausblick:}\\
Silberoxid (Ag$_2$O) ist das angelaufene Silber auf Schmuck oder Besteck.
Kann man diese Reaktion \textit{umkehren}? Kann man aus Silberoxid wieder reines Silber gewinnen?\\[0.3em]
$\rightarrow$ Das ist das Thema der nächsten Stunde: \textbf{Reduktion}
\end{infobox}

\end{document}
